\begin{afterword}
\summaryhead

The previous post showed that the cubes, whose first differences look patternless, have constant third differences. This post asks whether the real world, unlike closed mathematics, is simply messy. Its answer is no: physics has already found a simple, exact level beneath the surface, and the best guess is that ``reality is made of math.''

Why, then, can we not predict tigers and stock prices? The post counts three sources of uncertainty. We may not know the fundamental laws, but that matters little, since only particle accelerators reveal the gaps. We may lack the computing power to work out what known laws imply, as with protein folding: logical uncertainty. And we may not know where we are, like a little person on one of a sequence of cubes who knows all about the sequence but not which cube is theirs: indexical uncertainty, the subject of a book by Nick Bostrom.

Since biology's underlying order is already known and too deep to help, new physics will not make biology easy. Uncertainty is in the map, not the territory; the world is regular, and the Internet is ``a series of cubes.''

\responsehead

The post's three kinds of uncertainty are a useful way to sort the question, and its core claim, that failing to predict a system is no evidence that the system itself is messy, is sound. The indexical case goes back further than the book the post names: David Lewis's two gods (1979), who know every proposition true at their world but not which god they are. Its protein example has dated, since AlphaFold now predicts many structures from sequence, though by learning from known structures rather than from physics, so the point stands.

Two conclusions go further than the argument. The first is that ``reality is made of math.'' The paragraph shows that our best theories are mathematical and probably rest on deeper mathematics. That the world is a mathematical structure is a stronger claim, Max Tegmark's hypothesis, which is controversial and which the post does not argue for. Nothing else in the post depends on it.

The second is that, ``Because'' the master order is known, biology holds no further secret patterns that would make it easy. The reasons given, too little computing power and evolutionary accident, show that biology will not be derived from physics. They do not show that it lacks simple regularities at its own level. Thermodynamics found simple laws for systems whose molecular detail no one can compute. The post offers this as a guess, and the guess may be right, but its reasons do not reach it.

The quantum example takes many-worlds as what the theory ``quite definitely'' says, and the summary's ``deterministic'' depends on that reading. The conclusion, that certain quantum results cannot be predicted, holds on any standard reading, because the theory gives only probabilities. The author's case for many-worlds is discussed in the notes on ``Making Beliefs Pay Rent.''

In short: a useful sorting of uncertainty into three kinds, attached to two claims its argument does not reach, that reality ``is made of math'' and that biology holds no further simple patterns, and to many-worlds stated as what physics ``quite definitely'' says.
\end{afterword}
